% !TeX program = lualatex
% Compile twice: lualatex 222.tex
% All references and prose are embedded; no BibTeX database or external inputs.
\documentclass[11pt,a4paper]{article}
\usepackage[margin=24mm,headheight=14pt]{geometry}
\usepackage{fontspec}
\setmainfont{Latin Modern Roman}
\setsansfont{Latin Modern Sans}
\setmonofont{Latin Modern Mono}
\usepackage{amsmath,amssymb,mathtools}
\usepackage{unicode-math}
\setmathfont{Latin Modern Math}
\usepackage{microtype}
\usepackage{enumitem,xurl,needspace}
\usepackage[dvipsnames]{xcolor}
\usepackage{hyperref}
\hypersetup{unicode=true,colorlinks=true,linkcolor=MidnightBlue,urlcolor=MidnightBlue,
 pdftitle={Research Notebook: Section 222},pdfauthor={Research notebook},bookmarksnumbered=true}
\usepackage{bookmark}
\usepackage{tocloft}
\setlength{\cftsecnumwidth}{3em}
\cftsetpnumwidth{2.5em}
\cftsetrmarg{4em}
\usepackage{titlesec}
\titleformat{\section}{\Large\bfseries\raggedright}{\thesection}{0.7em}{}
\usepackage{fancyhdr}
\pagestyle{fancy}\fancyhf{}
\fancyhead[L]{\small\sffamily Open Problems | Research notebook}
\fancyhead[R]{\small 222 | 27 September 2026}
\fancyfoot[C]{\thepage}
\setlength{\parindent}{0pt}
\setlength{\parskip}{5pt plus 1pt}
\setlength{\emergencystretch}{3em}
\setcounter{tocdepth}{1}
\setcounter{secnumdepth}{1}
\setlist[itemize]{leftmargin=3em,itemsep=5pt,topsep=4pt,parsep=0pt}
\allowdisplaybreaks
\newcommand{\N}{\mathbb{N}}
\newcommand{\Z}{\mathbb{Z}}
\newcommand{\Q}{\mathbb{Q}}
\newcommand{\R}{\mathbb{R}}
\newcommand{\C}{\mathbb{C}}
\newcommand{\F}{\mathbb{F}}
\newcommand{\A}{\mathbb{A}}
\newcommand{\PP}{\mathbb P}
\newcommand{\E}{\mathbb{E}}
\newcommand{\eps}{\varepsilon}
\newcommand{\dd}{\,\mathrm d}
\newcommand{\sref}[2]{\hyperlink{source:#1:#2}{[S#2]}}
\newcommand{\eref}[2]{\hyperlink{extra:#1:#2}{[E#2]}}
% Turn the next switch off for a shorter reading copy. The quotations preserve
% every exactTarget field, including qualifications omitted from a short summary.
\newif\ifshowcatalogue
\showcataloguetrue
\newcommand{\cataloguescope}[1]{\ifshowcatalogue
 \paragraph{Exact scope in the supplied catalogue.}
 \begin{quote}\small #1\end{quote}\fi}

% Independently compilable extraction; original section text retained.
\numberwithin{equation}{section}
\begin{document}
\setcounter{section}{221}
% ================================================================
% problemId: problem.parshin-conjecture
% source releaseRank: 222; source releaseStatus: open
% exactTarget SHA-256: d1d0f0343f24a8059f9dd05e9b42236e096c17a9cbae7217fc3b29387a9a3c8b
\clearpage
\section{\texorpdfstring{Parshin Conjecture on Higher K-Theory over Finite Fields}{Parshin Conjecture on Higher K-Theory over Finite Fields}}\label{222}
\subsection{Definitions and mathematical statement}
For a smooth projective variety $X$ over a finite field $\F_q$, let $K_n(X)$ denote its higher algebraic $K$-group. An abelian group is torsion when every element is killed by some positive integer, with the integer allowed to depend on the element. Parshin's assertion is
\[
 K_n(X)\otimes_{\Z}\Q=0\quad\text{for all }n>0.
\]
This is equivalent to torsion, not necessarily finiteness of the integral group or a single bounded exponent. The group $K_0(X)$ is excluded, and properness and smoothness are part of the chosen scope. Passing to rational coefficients is the vanishing formulation, not an assertion of integral vanishing.
\par\smallskip\textit{Formulation and concept references:} \sref{222}{1}.
\cataloguescope{Prove that for every smooth projective variety X over a finite field and every integer n greater than zero, the higher algebraic K-group K\_n(X) is torsion.}
\subsection{Short English statement}
Are all positive-degree algebraic K-groups of smooth projective varieties over finite fields torsion?
% BEGIN RESEARCH ATTEMPT 222
\subsection{Research attempt}

\paragraph{Current frontier.}
The rational higher-Chow formulation and its relationship with
Weil--\'etale finite-generation conjectures are described by Geisser;
his Proposition~5.1 proves implications from stronger conjectures, not
an unconditional general vanishing theorem. \eref{222}{1}.
The 2024--2025 work of D\'eglise develops generic motives and motivic
cohomology of fields; results for those objects must not be identified
with the assertion for every smooth projective variety.
\eref{222}{4}, Introduction. The publisher page of \sref{222}{1} could
not be retrieved in this session. The supplied statement is
self-contained, and no unnamed theorem from that inaccessible paper is
used below. The frontier check therefore has that explicit limitation.

\paragraph{Attack route.}
The three routes tested are finite-field descent, motivic decomposition,
and comparison with $\ell$-adic cohomology. Descent works rationally but
does not find a favorable extension. A decomposition of the diagonal
through motives with known vanishing would suffice, provided its error
is genuinely nilpotent. The $\ell$-adic target can be shown to vanish by
weights; this makes injectivity of comparison especially attractive,
but also exposes the invisible-divisibility obstruction. The attempt
makes these distinctions explicit rather than replacing torsion by
integral finiteness.

\paragraph{Attempt.}
\textbf{1. Exact descent and projective-bundle subcases.}
Let $f:X_{\F_{q^a}}\to X$ be the constant-field extension of degree $a$.
The transfer and pullback satisfy
\[
 f_*f^*(z)=a z\qquad(z\in K_n(X)).
\]
Indeed $f_*\mathcal O_{X_{\F_{q^a}}}$ is a free $\mathcal O_X$-module
of rank $a$, and the projection formula gives the identity. Hence
$f^*:K_n(X)\otimes\Q\to K_n(X_{\F_{q^a}})\otimes\Q$ is injective.
Proving rational vanishing after a finite extension would therefore
suffice; no uniform extension degree is required by this deduction.
Transfers and the projection formula are recorded in \eref{222}{2},
Section~V.3.

The projective-bundle formula also gives
\[
 K_n(\PP^r_X)\cong\bigoplus_{j=0}^r K_n(X).
\]
In particular a tower of projective bundles over a finite field has the
asserted positive-degree rational vanishing, using Quillen's calculation
$K_{2i}(\F_q)=0$ and
$K_{2i-1}(\F_q)\cong\Z/(q^i-1)$ for $i\ge1$.
\eref{222}{2}, Projective Bundle Theorem~V.1.5 and the finite-field
calculation in Chapter~IV. These familiar subcases are a useful test,
not a proof that all smooth projective varieties become such towers.

\textbf{2. A nilpotent-diagonal criterion.}
Write $CH^r(X,n)_{\Q}$ for rational higher Chow groups. The rational
Chern-character/Adams decomposition identifies positive $K_n(X)_{\Q}$
with the sum of these groups over $r$. Thus their vanishing for every
$r$ is the relevant stronger-looking but equivalent graded target.
\eref{222}{1}, Sections~1 and~3.2; \eref{222}{4}, motivic-cohomology
conventions.

Use rational Chow correspondences, with Tate twists chosen so that the
higher-Chow functor on a summand $\mathbb L^j$ is
$CH^{r-j}(\operatorname{Spec}\F_q,n)_{\Q}$.
Suppose there are morphisms of rational Chow motives
\[
 M(X)\xrightarrow{u}T\xrightarrow{v}M(X),\qquad
 T=\bigoplus_{j\in J}\mathbb L^j,
\]
for a finite multiset $J$, and that
\begin{equation}\label{222:nilpotent}
 e=1_{M(X)}-v u,\qquad e^N=0
\end{equation}
for some finite $N$. Correspondences act functorially on higher Chow
groups by pullback, intersection and proper pushforward; their
composition is the usual correspondence composition. For $n>0$, the
higher-Chow groups of $T$ vanish, by the finite-field calculation and
its rational grading. Therefore the action of $vu$ on
$CH^r(X,n)_{\Q}$ is zero. Equation \eqref{222:nilpotent} says that the
identity action on this group is nilpotent. Taking its $N$th power
proves that the group is zero. Thus \eqref{222:nilpotent} is a proved
sufficient condition for Parshin's assertion for $X$.

There is also a useful exact correction of the motive maps:
\[
 (1+e+\cdots+e^{N-1})vu=1-e^N=1.
\]
It makes $M(X)$ a retract of $T$. This shows how strong the proposed
condition really is. A general variety cannot be assumed to have such
a Tate domination; an elliptic curve, for instance, has nonzero odd
$\ell$-adic cohomology, whereas finite direct sums of pure Tate motives
do not. This does not contradict Parshin for curves. It shows that a
universal proof cannot obtain \eqref{222:nilpotent} with only Tate
intermediates for every $X$. A more flexible route would have to use
other motives with independently established positive-degree vanishing.

\textbf{3. The comparison target vanishes, but comparison may forget everything.}
Fix $\ell\ne\operatorname{char}\F_q$. For $n>0$ and an integer $r$,
put $m=2r-n$. If $m<0$, the usual $\ell$-adic group in that degree is
zero. If $m\ge0$, the Hochschild--Serre sequence over a finite field
expresses $H^m_{\mathrm{et}}(X,\Q_\ell(r))$ by the coinvariants of
$H^{m-1}_{\mathrm{et}}(\overline X,\Q_\ell(r))$ and the invariants of
$H^m_{\mathrm{et}}(\overline X,\Q_\ell(r))$.
Smoothness and properness give pure Frobenius weights. The eigenvalue
absolute values on these two spaces are respectively
\[
 q^{(-n-1)/2}\quad\hbox{and}\quad q^{-n/2}.
\]
Neither is one. Thus Frobenius minus identity is invertible, and
\begin{equation}\label{222:etale}
 H^{2r-n}_{\mathrm{et}}(X,\Q_\ell(r))=0\qquad(n>0).
\end{equation}
This calculation uses the established weight theorem, not Parshin.
\eref{222}{3}, smooth projective case of the Weil conjectures.

Consequently an injective rational cycle/regulator map
\[
 CH^r(X,n)\otimes\Q_\ell\longrightarrow
 H^{2r-n}_{\mathrm{et}}(X,\Q_\ell(r))
\]
would settle the desired vanishing. But in this range such an injectivity
statement is precisely the hard assertion; its target is already zero.
Finite-coefficient comparison cannot be used to manufacture it. The
abelian group $V=\Q$ has
\[
 V/\ell^aV=0,\qquad V[\ell^a]=0\quad(a\ge1),
 \qquad V\otimes\Q=\Q\ne0.
\]
Thus even all finite quotients and all finite torsion tests can be blind
to a rational summand. Completion need not inject into or preserve the
original group. This is a counterexample to an inference about abelian
groups, not a construction of such a summand inside $K_n(X)$.

An alternative stronger assumption, finite generation in the appropriate
Weil--\'etale theory, is explicitly known to imply Parshin in
\eref{222}{1}, Proposition~5.1. Using that implication is legitimate
conditionally; declaring that finite generation established for every
$X$ would simply assume another unresolved conjecture.

\paragraph{Stress test.}
The transfer is applied to a constant-field extension, whose pushforward
of the structure sheaf really is free of rank $a$. A general generically
finite map cannot be substituted without checking the relevant
push-pull identity. Nilpotence in \eqref{222:nilpotent} is an equality
of rational correspondences, not merely a vanishing of their numerical
or $\ell$-adic realizations. Its exponent is finite; no infinite-series
convergence is used.

The weight calculation depends on both smoothness and properness. It
does not predict the same vanishing for arbitrary open varieties.
The group $K_0$ is not included: its cycle classes lie in weight zero
in the relevant degree. No uniform exponent killing integral $K_n$, or
integral finiteness, is inferred from rational vanishing.

\paragraph{Outcome.}
\textbf{Outcome: Conditional reduction.}

The attempt proves a nilpotent-correspondence sufficient condition,
checks finite-extension descent, and identifies exactly why vanishing
$\ell$-adic comparison targets do not themselves prove the conjecture.
The projective-bundle consequences are established subcases, not a new
unrestricted result or a claim of priority.

What remains is a decomposition through sufficiently general motives
with known vanishing, or a theorem excluding the rational classes that
finite-coefficient and completed realizations cannot see. Neither has
been supplied for an arbitrary smooth projective $X$.
\needspace{26\baselineskip}
% END RESEARCH ATTEMPT 222
\subsection{Sources}
\begin{itemize}\raggedright
\item[\textnormal{[S1]}]\hypertarget{source:222:1}{} On Parshin's conjecture.
\url{https://www.sciencedirect.com/science/article/pii/S0022404923002001}.
{\small\textit{Catalogue formulation source.}}
% BEGIN ADDED REFERENCES 222
\item[\textnormal{[E1]}]\hypertarget{extra:222:1}{}
T. Geisser, \emph{Finite generation conjectures for cohomology over finite
fields}, arXiv:1103.5544v1, 29 March 2011, Sections~2.2, 3.2 and
Proposition~5.1. \url{https://arxiv.org/html/1103.5544v1}.
\item[\textnormal{[E2]}]\hypertarget{extra:222:2}{}
C. A. Weibel, \emph{The K-book: An Introduction to Algebraic K-theory},
Graduate Studies in Mathematics 145, AMS, 2013; author chapters IV
(finite fields) and V, Theorem~1.5 and Section~3 (projective bundles and
transfers). \url{https://sites.math.rutgers.edu/~weibel/Kbook/Kbook.IV.pdf};
\url{https://sites.math.rutgers.edu/~weibel/Kbook/Kbook.V.pdf}.
\item[\textnormal{[E3]}]\hypertarget{extra:222:3}{}
P. Deligne, \emph{La conjecture de Weil. I}, Publications Math\'ematiques
de l'IH\'ES \textbf{43} (1974), 273--307, the Frobenius weight theorem
for smooth projective varieties.
\url{https://www.numdam.org/item/PMIHES_1974__43__273_0/}.
\item[\textnormal{[E4]}]\hypertarget{extra:222:4}{}
F. D\'eglise, \emph{Generic motives and motivic cohomology of fields},
arXiv:2408.06233v2, 11 January 2025, Introduction and motivic-cohomology
conventions. \url{https://arxiv.org/html/2408.06233v2}.
All four added sources consulted 27 September 2026. Conditional statements
in them are not treated as unconditional theorems.
% END ADDED REFERENCES 222
\end{itemize}

\end{document}
